% NB14a -- Arquitectura del PT6A-61 (flujo inverso, dos ejes) con la numeracion de estaciones de P&WC.
% El aire entra por detras (derecha), avanza hacia la helice (izquierda): estaciones 1 -> 7.
\documentclass[tikz,border=4mm]{standalone}
\usepackage{amsmath,amssymb}
\usetikzlibrary{arrows.meta,calc,positioning,decorations.pathmorphing}
\begin{document}
% Estilo comun de los diagramas TikZ del modulo VCR (NB10b) -- paleta del proyecto
\definecolor{manivela}{HTML}{2A78D6}
\definecolor{inferior}{HTML}{4A3AA7}
\definecolor{control}{HTML}{1BAF7A}
\definecolor{superior}{HTML}{EB6834}
\definecolor{tinta}{HTML}{3A3A38}
\definecolor{metal}{HTML}{C3C2B7}
\definecolor{perdida}{HTML}{E34948}
\definecolor{fuerza}{HTML}{C0392B}
\tikzset{
  barra/.style={line width=3.2pt, line cap=round},
  pin/.style={circle, fill=white, draw=tinta, line width=0.9pt, inner sep=1.6pt},
  fijo/.style={circle, fill=tinta, inner sep=2.6pt},
  vec/.style={-{Stealth[length=2.6mm]}, line width=0.9pt},
  F/.style={-{Stealth[length=3mm]}, line width=1.2pt, fuerza},
  cota/.style={{Stealth[length=1.8mm]}-{Stealth[length=1.8mm]}, line width=0.5pt, tinta},
  lbl/.style={font=\small, text=tinta},
}
\newcommand{\tierra}[1]{\fill[pattern=north east lines, pattern color=tinta] ($(#1)+(-5,-3)$) rectangle ($(#1)+(5,-6)$);
  \draw[tinta, line width=0.8pt] ($(#1)+(-5,-3)$) -- ($(#1)+(5,-3)$);}

\definecolor{frio}{HTML}{2A78D6}
\definecolor{gg}{HTML}{4A3AA7}
\definecolor{caliente}{HTML}{E34948}
\definecolor{pt}{HTML}{EB6834}
\definecolor{helice}{HTML}{1BAF7A}
\begin{tikzpicture}[x=1cm,y=1cm,
  bloque/.style={draw=tinta, line width=0.9pt, rounded corners=2pt, align=center, font=\scriptsize, minimum height=16mm},
  est/.style={font=\footnotesize\bfseries, text=tinta},
  flujo/.style={-{Stealth[length=2.6mm]}, line width=1.4pt},
  tag/.style={font=\scriptsize\itshape, text=tinta!75}]
  % eje central
  \draw[tinta!50, dash dot] (-1.2,0) -- (17.2,0);
  % helice
  \draw[helice, line width=3pt] (-0.6,-2.6) -- (-0.6,2.6);
  \fill[helice] (-0.6,0) circle (2.5mm);
  \node[tag, helice!70!black] at (-0.6,3.0) {h\'elice McCauley, 4 palas, $R=1.193$ m};
  % reductor
  \node[bloque, fill=metal!40, minimum width=17mm, minimum height=22mm] (red) at (0.9,0) {reductor\\planetario\\2 etapas\\$15.07{:}1$};
  % escape
  \node[bloque, fill=pt!15, minimum width=12mm] (esc) at (2.9,0) {escape\\(2 toberas\\laterales)};
  % turbina de potencia
  \node[bloque, fill=pt!35, minimum width=19mm] (tp) at (4.85,0) {turbina\\libre (TP)\\2 etapas};
  % turbina del compresor
  \node[bloque, fill=gg!25, minimum width=14mm] (tc) at (6.95,0) {turbina\\del\\compresor\\(TC)};
  % camara de flujo inverso
  \node[bloque, fill=caliente!18, minimum width=30mm, minimum height=12mm] (cc) at (9.35,1.55) {c\'amara anular de\\flujo inverso, $V=0.0288$ m$^3$};
  % compresor centrifugo
  \node[bloque, fill=frio!35, minimum width=14mm] (cen) at (11.75,0) {compresor\\centr\'ifugo};
  % compresor axial
  \node[bloque, fill=frio!20, minimum width=26mm] (ax) at (14.0,0) {compresor axial\\3 etapas};
  % admision
  \node[bloque, fill=frio!8, minimum width=12mm] (adm) at (16.15,0) {admisi\'on\\(malla)};
  % ejes concentricos
  \draw[gg, line width=2.4pt] (6.95,-1.25) -- (14.0,-1.25);
  \node[tag, gg] at (10.5,-1.6) {eje del generador de gas $N_g$ (100\,\% = 37\,468 rpm)};
  \draw[pt, line width=2.4pt] (0.9,-2.05) -- (4.85,-2.05);
  \node[tag, pt!80!black] at (2.9,-2.4) {eje de potencia $N_p$ (30\,145 rpm $\leftrightarrow$ 2000 rpm h\'elice)};
  \draw[gg, line width=1pt] (6.95,-0.8) -- (6.95,-1.25); \draw[gg, line width=1pt] (14.0,-0.8) -- (14.0,-1.25); \draw[gg, line width=1pt] (11.75,-0.8) -- (11.75,-1.25);
  \draw[pt, line width=1pt] (4.85,-0.8) -- (4.85,-2.05); \draw[pt, line width=1pt] (0.9,-1.1) -- (0.9,-2.05);
  % flujo de aire / gas
  \draw[flujo, frio] (17.4,0.55) -- (16.85,0.55);
  \draw[flujo, frio] (12.3,0.95) .. controls (12.6,2.3) and (11.6,2.4) .. (10.9,2.25);
  \draw[flujo, caliente] (7.8,1.55) .. controls (7.2,1.5) and (7.2,0.95) .. (7.55,0.95);
  \draw[flujo, caliente] (2.9,0.85) -- (2.9,2.0);
  \node[tag] at (17.0,1.0) {aire};
  \node[tag, caliente] at (2.9,2.3) {gases};
  % estaciones
  \foreach \x/\s in {17.0/1, 16.6/2, 14.9/2.5, 13.1/2.8, 11.1/3, 7.75/4, 6.1/5, 3.8/6, 2.35/7}{
     \node[est, fill=white, inner sep=1pt, circle, draw=tinta, line width=0.5pt] at (\x,-3.1) {\s};
     \draw[tinta!60] (\x,-2.9) -- (\x,-0.82);
  }
  \node[tag, anchor=west] at (-1.2,-3.6) {Estaciones de P\&WC (carta de despegue, d\'ia est\'andar, nivel del mar): 1 ambiente, 2 entrada al compresor, 2.5/2.8 entre etapas axiales, 3 descarga del compresor,};
  \node[tag, anchor=west] at (-1.2,-4.0) {4 entrada a la TC, 5 entre turbinas (se mide la ITT), 6 salida de la TP, 7 escape.};
\end{tikzpicture}
\end{document}
